\hdr{§11}{考场执行：150 分钟怎么花}

\textbf{这一章明早出门前读一遍即可。}知识已经定了，能再提分的是\emph{顺序、时间、写法}这三件事。

\needspace{11\baselineskip}
\subhdr{11.1 时间预算}

\begin{center}
\small
\tabtop
\begin{tabular}{@{}p{2.4cm}p{2.2cm}p{8.4cm}@{}}
\thead{时段} & \thead{做什么} & \thead{要点}\\
\hline
0--5 min & 通览全卷 & 数清题号，标出“一眼会”“要想想”“完全不会”三类；不动笔\\
5--35 min & 填空题 5 道 & 每题 $\le 5$ 分钟；卡住就跳过，绝不纠缠\\
35--55 min & 最容易的两道大题 & 先拿稳的分，建立节奏\\
55--120 min & 剩下三道大题 & 每题 18--20 分钟；到点无论写没写完都换题\\
120--145 min & 回头补 & 补填空、补没写完的步骤\\
145--150 min & 检查 & 只查符号、上下限、常数、方向\\
\end{tabular}
\tabbot
\end{center}

\warnbox{\textbf{一条铁律}：单题超过 12 分钟没有实质进展，立刻换题。竞赛卷的分数分布是“每道题都有 3--5 分的入门步骤”，把时间花在已经卡住的题上，等于放弃了后面整道题的入门分。}

\subhdr{11.2 解答题的标准写法（阅卷按步给分）}

无论会不会做完，每题都按这个骨架写，能拿到过程分：

\begin{enumerate}[leftmargin=1.8em,itemsep=2pt]
\item \textbf{抄条件}：把题目给的条件用数学语言重写一遍（“$f$ 在 $[0,1]$ 连续，在 $(0,1)$ 可导，$f(0)=0$”）。这一步就有分。
\item \textbf{写目标}：把要证／要求的东西写成一行式子。
\item \textbf{写构造}：辅助函数、坐标变换、补的线或面，明确写出来并说明为什么这样构造（“令 $F(x)=e^{-x}f(x)$，则 $F(0)=F(1)=0$”）。
\item \textbf{写依据}：点名用哪个定理，并核对其条件（“由罗尔定理，存在 $\xi$……”）。
\item \textbf{写结论}：把结果整理成题目要求的形式，加一句“故得证”。
\end{enumerate}

\textbf{不会做时的三步抢救}：
\begin{itemize}[leftmargin=1.6em,itemsep=1pt]
\item 写出相关定义（导数定义、可微定义、格林公式的条件），常能拿到 1--2 分；
\item 写出特殊情形（$n=1$、$f$ 为常数、区域为圆）的结果，可能触发思路；
\item 把能算的中间量算出来（如 $Q_x-P_y$、特征根、收敛半径），即使最后没积完。
\end{itemize}

\subhdr{11.3 最后 15 分钟检查什么}

按这个顺序查，每项 1--2 分钟：

\begin{enumerate}[leftmargin=1.8em,itemsep=1pt]
\item \textbf{符号}：格林公式方向、高斯公式内外侧、第二类曲面积分的正负号。
\item \textbf{上下限}：换元后是否换了限；极坐标 $r$ 的上限是否含 $\theta$；换序后是否写反。
\item \textbf{常数}：不定积分 $+C$；微分方程 $C_1,C_2$；幂级数求和时代入的 $x$ 是否在收敛域内。
\item \textbf{特解}：二阶非齐次的 $x^k$ 是否漏掉。
\item \textbf{答案形式}：是否要求“通解”却只给了“特解”；是否要求“最大值”却给了驻点。
\end{enumerate}

\subhdr{11.4 十二个高频失分点（一页背完）}

\begin{center}
\small
\tabtop
\begin{tabular}{@{}p{0.7cm}p{6.0cm}p{6.0cm}@{}}
\thead{\#} & \thead{错误} & \thead{正确}\\
\hline
1 & 等价无穷小在加减法里替换 & 先提公因子，或改用泰勒\\
2 & 洛必达不验证 $0/0$、$\infty/\infty$ & 先看型，再求导\\
3 & 泰勒展开阶数不够，出现 $0/0$ & 展开到分母同阶\\
4 & 罗尔定理没验证 $F(a)=F(b)$ & 构造后先算两端值\\
5 & 积分换元不换限 & 换元同时换上下限\\
6 & 极坐标忘乘 $r$ & $\dd x\dd y=r\dd r\dd\theta$\\
7 & 格林公式方向取反 & 逆时针为正\\
8 & 高斯公式用内侧 & 外侧为正\\
9 & 幂级数不查端点 & $x=\pm R$ 单独判\\
10 & ODE 特解漏 $x^k$ & $k$ 取 $\lambda$ 的重数\\
11 & 条件极值漏约束方程 & $L_x=L_y=L_z=0$ 加 $g=0$\\
12 & 证明题只写结论不写依据 & 点名定理 + 核条件\\
\end{tabular}
\tabbot
\end{center}

\callout{\textbf{明早出发前 10 分钟只看这三样}：附录 A 的速查表、“二阶常系数特解设法”表、上面这 12 条。看完就合上，不要再看新题——新题只会增加焦虑，不会增加分数。}
